\chapter{INTRODUCTION}
\label{ch1}

Write your introductory chapter here. For example, your literature
overview could go here. Sometimes You will want a section to go within your chapter.

\section{Section 1}

Section 1 goes here.

Here is an example of a table.  These are the UD holidays for AY 2017-2018.

\begin{table}
    \begin{center}
        \caption{University of Dayton Holidays for Academic Year 2017-2018}
        \label{tab:class}
        \begin{tabular}{||l|c|c|r||}
            \hline
            Holdiay & Month & Day & Year\\
            \hline
            \hline
            Labor Day & September & 4 & 2017 \\
            \hline
            Mid-Terms Break & October & 5-6 & 2017 \\
            \hline
            Thanksgiving & November & 22-24 & 2017 \\
            \hline
            Christmas Break & December & 17 & 2017 \\
            \hline
            Christmas Break Ends & January & 16 & 2018 \\
            \hline
            Spring Break & February & 28 & 2018 \\
            \hline
            Spring Break Ends & March & 5 & 2018 \\
            \hline
            Easter Recess & March & 28 & 2018 \\
            \hline
            Easter Recess Ends & April & 3 & 2018 \\
            \hline
            Stander Symposium & April & 18 & 2018 \\
            \hline
            \multicolumn{4}{} \\
        \end{tabular}
     \end{center}
     \alt{This is a table of University Holidays}
\end{table}

We can reference the University holidays in Table \ref{tab:class}.

Here is a second table showing the UD Men's 2018-2019 basketball schedule

\section{Section 2}

\begin{table}
    \begin{center}
        \caption{University of Dayton Men's Basketball Schedule 2018-2019}
        \label{tab:bball}
        \begin{tabular}{||l|c||}
            \hline
            Date & Opponent \\
            \hline \hline
            November 2 & Capital \\
            \hline
            November 7 & North Florida \\
            \hline
            November 10 & Coppin State \\
            \hline
            November 16 & Purdue Fort Wayne \\
            \hline
            November 21 & Butler \\
            \hline
            November 22 & Virginia \\
            \hline
            November 23 & Oklahoma \\
            \hline \hline
            December 4 & Detroit Mercy \\
            \hline
            December 8 & Auburn \\
            \hline
            December 16 & Tulsa \\
            \hline
            December 19 & Western Michigan \\
            \hline
            December 22 & Presbyterian \\
            \hline
            December 29 & Georgia Southern \\
            \hline \hline
            January 6 & Richmond \\
            \hline
            January 9 & George Washington \\
            \hline
            January 13 & Massachusetts \\
            \hline
            January 16 & VCU \\
            \hline
            January 19 & St. Bonaventure \\
            \hline
            January 23 & George Mason \\
            \hline
            January 26 & Fordham \\
            \hline
            January 29 & Saint Joseph's \\
            \hline \hline
            February 2 & Duquesne \\
            \hline
            February 5 & Saint Louis \\
            \hline
            February 9 & Rhode Island \\
            \hline
            February 16 & VCU \\
            \hline
            February 19 & Davidson \\
            \hline
            February 23 & Saint Louis \\
            \hline
            February 26 & Massachusetts \\
            \hline \hline
            March 1 & Rhode Island \\
            \hline
            March 6 & La Salle \\
            \hline
            March 9 & Douquesne \\
            \hline
            \multicolumn{2}{} \\
        \end{tabular}
     \end{center}
     \alt{This is the 2017-2018 UD Men's Basketball schedule}
\end{table}

Again, we can reference basketball schedule as Table \ref{tab:bball}.

Section 2 goes here.

An example of an equation is 
\begin{equation}
    x[n] = \left\{ \begin{array}{ll}
        \infty, & {\tt t=0} \\
        0, & {\tt else}
    \end{array}\right..
    \label{eqn:delta}
\end{equation}

Now, we can reference this equation as Equation \ref{eqn:delta}.

Here is an example of a citation.  In \cite{JMRK2002}, the authors discuss adaptive control using on-line
function approximation.
